% year: תשפ"ג
% type: מבחן
% semester: א
% moed: א
% number: 6ב
% points: 10
% categories: improper-integrals
% source: exams/מבחנים/תשפג/מועד א תשפג.pdf

\begin{question}
חשבו את האינטגרל הלא-אמיתי $\displaystyle\int_0^{\infty}e^{-2x}\sin x\,dx$ או הראו שאינו מתכנס.
\end{question}

\begin{hint}
חשבו את האינטגרל עד $R$ בעזרת אינטגרציה בחלקים פעמיים (האינטגרל "חוזר על עצמו"), ואז העבירו $R\to\infty$.
\end{hint}

\begin{solution}
\textbf{התכנסות:} $|e^{-2x}\sin x|\le e^{-2x}$ ו-$\int_0^\infty e^{-2x}dx=\frac12$ מתכנס, ולכן לפי \textbf{מבחן ההשוואה} האינטגרל מתכנס (בהחלט).

\textbf{חישוב:} נסמן $I(R)=\int_0^R e^{-2x}\sin x\,dx$ ונמצא פונקציה קדומה $J=\int e^{-2x}\sin x\,dx$ באינטגרציה בחלקים פעמיים ($u=e^{-2x}$):
\begin{align*}
J&=-e^{-2x}\cos x-\int 2e^{-2x}\cos x\,dx
 =-e^{-2x}\cos x-2\left(e^{-2x}\sin x+\int 2e^{-2x}\sin x\,dx\right)\\
 &=-e^{-2x}\cos x-2e^{-2x}\sin x-4J.
\end{align*}
לכן $5J=-e^{-2x}(\cos x+2\sin x)$, כלומר
\[ \int e^{-2x}\sin x\,dx=-\frac{e^{-2x}(\cos x+2\sin x)}{5}+C. \]
מכאן
\[ I(R)=\frac{1}{5}-\frac{e^{-2R}(\cos R+2\sin R)}{5}. \]
כאשר $R\to\infty$: $|\cos R+2\sin R|\le 3$ ו-$e^{-2R}\to 0$, ולכן (חסומה כפול שואפת לאפס) האיבר השני שואף ל-$0$. לכן
\[ \int_0^\infty e^{-2x}\sin x\,dx=\frac15. \]
\end{solution}
